% deck_headline.tex -- the causal question, the evidence and what we can claim (advisor meeting). Built 2026-11-23.
% Exhibits: generated/headline__*.tex and generated/assets/* (script4o-headline-deck.py, script4e-tables.py, script4f-figures.py); no number is typed here
% except in the frame text that quotes a table beside it. Compile: xelatex deck_headline.tex (Gotham theme).
\documentclass[aspectratio=169]{beamer}
\usetheme{gotham}
\gothamset{colorset=anthracite}
\gothamset{background=transparent}
\gothamset{progressbar position=foot, progressbar style=rounded box,
           progressbar advancement=brlt, numbering=totalframenumber}
\boolfalse{sectionContent}
\setbeamerfont{itemize/enumerate body}{size=\small}
\setbeamerfont{itemize/enumerate subbody}{size=\footnotesize}

\usepackage{booktabs,array,adjustbox,multirow,graphicx,amsmath,amssymb}

\newcommand{\fig}[2][0.42]{\begin{center}\includegraphics[width=\textwidth,height=#1\textheight,keepaspectratio]{#2}\end{center}}
\newcommand{\figcols}[2]{\begin{columns}[T]\begin{column}{0.64\textwidth}\includegraphics[width=\textwidth]{#1}\end{column}\begin{column}{0.33\textwidth}\setbeamerfont{itemize/enumerate body}{size=\scriptsize}#2\end{column}\end{columns}}
\newcommand{\src}[1]{{\tiny #1\par}}

\defbeamertemplate{title page}{}{\centering{\huge \inserttitle}\\[0.6em]{\large \insertsubtitle}\\[1em]{\large \insertauthor}\\[1em]{\insertdate}}
\title{Do large dairies raise poor mental health days?}
\subtitle{The causal claim, the evidence, and what we can say}
\author{Allegra Saggese for Professors Galina Hale, Hanna Kim}
\date{\today}

\begin{document}

\begin{frame}\titlepage\end{frame}

\begin{frame}{Agenda}
\tableofcontents
\end{frame}

% =====================================================================
\section{The causal claim}

\begin{frame}{The claim we want to test}
\begin{itemize}\small\setlength\itemsep{4pt}
  \item[] \textbf{Claim.} When a rural U.S.\ county gets its first large dairy operation, or its first expansion of large dairies, residents report \emph{more poor mental health days}.
  \item[] \textbf{Estimand.} The average effect on counties that were treated, in days per month, against counties that were not yet treated (or never treated).
  \item[] \textbf{Identifying assumption.} Without the dairy, treated and comparison counties would have moved together (parallel trends), once we hold fixed eight pre-committed county characteristics.
  \item[] \textbf{Possible channels, not tested.} Environmental stress, dependence on one industry, farm consolidation.
  \item[] \textbf{Second claim, tested separately.} Slaughterhouse plants raise distress via workers.
\end{itemize}
\end{frame}

\begin{frame}{What would make us believe it, and what would refute it}
\begin{columns}[T]
\begin{column}{0.48\textwidth}
\textbf{We would believe it if}
\begin{itemize}\footnotesize\setlength\itemsep{2pt}
  \item the effect shows up under estimators built for staggered timing, not only in two-way fixed effects
  \item treated and comparison counties were moving together \emph{before} the dairy arrived
  \item it survives the number of tests we ran
  \item it is specific to large dairies, not beef, hogs or chickens
  \item it survives spatial checks and the 2016 change in how the outcome is measured
\end{itemize}
\end{column}
\begin{column}{0.48\textwidth}
\textbf{We would reject it if}
\begin{itemize}\footnotesize\setlength\itemsep{2pt}
  \item the estimate depends on one estimator
  \item pre-period estimates are already different from zero
  \item the result does not survive correction across the outcomes and treatments tested
  \item other animals show the same pattern
  \item it exists on only one side of the 2016 data break
\end{itemize}
\end{column}
\end{columns}
\end{frame}

% =====================================================================
\section{What data we have}

\begin{frame}{County-year panel: outcomes and samples}
\input{generated/headline__data_outcomes.tex}
\end{frame}

\begin{frame}{Treatment events: dairies and plants}
\input{generated/headline__data_events.tex}
\end{frame}

\begin{frame}{Where the first expansions happen}
\fig[0.55]{generated/assets/expansion_cohort_year.png}
\src{First year a county's count of large dairies rises at a census wave (2012, 2017 or 2022 inside the Poor MH Days window). Source: panel 2026-10-25; script1d-maps.py.}
\end{frame}

\begin{frame}{What the data cannot see}
\begin{itemize}\small\setlength\itemsep{4pt}
  \item \textbf{Dairy timing is coarse.} Operation counts come from the Census of Agriculture, every five years. A county ``first expands'' in a census year even if the expansion happened in the five years before it.
  \item \textbf{The outcome changes method.} The county mental-health series is survey-based before 2016 and CDC model-based after; there is a second change in 2021.
  \item \textbf{Slaughterhouses are observed only 2017--2023,} so most plant-entry cohorts have one to three pre-years.
  \item \textbf{Crime is arrests, not offences,} and its pre-trends fail; it is not part of the headline.
  \item \textbf{Rural counties only,} so a plant in an urban county of a commuting zone is invisible.
\end{itemize}
\end{frame}

% =====================================================================
\section{How we test the claim}

\begin{frame}{Research design: staggered difference-in-differences}
\footnotesize
\textbf{Two-way fixed effects (TWFE):} $y_{it}=\alpha_i+\delta_t+\beta D_{it}+X_{it}'\gamma+\varepsilon_{it}$, county ($\alpha_i$) and year ($\delta_t$) effects, standard errors clustered by state.\par
\vspace{4pt}
\textbf{Why not stop there:} counties are treated at different times, so TWFE also compares newly treated counties with counties treated earlier.\par
\vspace{4pt}
\textbf{Callaway--Sant'Anna (CS):} for each first-treated cohort $g$ and year $t$,
\[ ATT(g,t)=E[\,Y_t-Y_{g-1}\mid G=g\,]-E[\,Y_t-Y_{g-1}\mid \text{not yet treated}\,], \]
adjusted for the eight controls, averaged over cohorts and years; standard errors by bootstrap over states.\par
\vspace{4pt}
\textbf{Treatments:} E2 first expansion (first census wave with a rise in large dairies; every county) and E3 first entry (first large dairy in a county with none; clean controls).
\end{frame}

\begin{frame}{The test plan: six tests, each answered on a later frame}
\begin{center}\adjustbox{max width=\textwidth}{\footnotesize
\begin{tabular}{@{}lll@{}}
\toprule
Test & What it checks & Reported on\\
\midrule
Estimator & Does the effect survive estimators built for staggered timing? & ``Estimators disagree''\\
Pre-trends & Were treated and comparison counties moving together before? & Two pre-trend frames\\
Many tests & Does it survive the family of outcomes and treatments tested? & ``Multiple testing''\\
Adjustments & Controls, region-by-year effects, spatial errors, dropping nearby counties & ``What holds up''\\
Specificity & Is it large dairies only? & ``What holds up''\\
Measurement & Same on both sides of the 2016 data break? & ``2016 data break''\\
Slaughterhouses & Large-plant entry, zone exposure, closures & Last results frames\\
\bottomrule
\end{tabular}}
\end{center}
\end{frame}

\begin{frame}{Slaughterhouse designs, with the rules fixed before results}
\begin{itemize}\footnotesize\setlength\itemsep{4pt}
  \item \textbf{Large plant:} an FSIS establishment of size code 4 or 5 (large, very large), 2017--2023.
  \item \textbf{County entry and exit:} counties without (with) a large plant in 2017; treated from the first year one appears (disappears). \emph{D-63}
  \item \textbf{Commuting-zone exposure:} treat the zone, so workers living next door count; version A all counties in the zone, version B only counties that never host a plant. \emph{D-64}
  \item \textbf{Closures:} consolidation means some areas lost plants and jobs. We either drop closure areas from the comparison or control for a closure indicator. \emph{D-65, D-66}
  \item \textbf{Many plants versus none:} counties with two or more large plants against counties with none; read only if the pre-period test is flat. \emph{D-65}
\end{itemize}
\end{frame}

% =====================================================================
\section{Interpreting the results}

\begin{frame}{Estimators disagree: first expansion}
\figcols{generated/assets/forest_triangulation_first_expansion__main.pdf}{
\begin{itemize}\setlength\itemsep{4pt}
  \item Poor MH Days (left panel): TWFE 0.104 days (0.096 in this sample); Callaway--Sant'Anna 0.037, about a third of it, not significant.
  \item Sun--Abraham and de Chaisemartin--D'Haultfoeuille agree with Callaway--Sant'Anna; imputation and stacked DiD sit near TWFE.
\end{itemize}}
\end{frame}

\begin{frame}{Pre-trends (1): TWFE event study}
\figcols{generated/assets/event_study_twfe__main.pdf}{
\begin{itemize}\setlength\itemsep{4pt}
  \item Poor MH Days: the joint test of the seven leads does not reject (p = 0.43 for E2, 0.61 for E3).
  \item Pre-period points drift upward toward the event; the joint test does not call it a trend.
\end{itemize}}
\end{frame}

\begin{frame}{Pre-trends (2): Callaway--Sant'Anna event study}
\figcols{generated/assets/event_study_cs__main.pdf}{
\begin{itemize}\setlength\itemsep{4pt}
  \item Poor MH Days: 6 of 11 pre-period estimates are individually significant (1 of 11 without covariates).
  \item Possible reason, not tested: a census-wave ``first expansion'' may be up to five years old, so some pre-periods are post-expansion years.
\end{itemize}}
\end{frame}

\begin{frame}{Multiple testing: one result survives, and it fails its lead test}
\input{generated/headline__multiple.tex}
\begin{itemize}\footnotesize\setlength\itemsep{1pt}
  \item The log-count treatment (T4) is the only Poor MH Days result with a Holm $p$ below 0.05.
  \item Its lead test fails ($p$ = 0.027), so the one test that survives correction is the one with a visible pre-trend.
\end{itemize}
\end{frame}

\begin{frame}{What holds up (1): region-by-year effects, spatial errors, nearby counties}
\input{generated/headline__robust_adj.tex}
\begin{itemize}\footnotesize\setlength\itemsep{1pt}
  \item TWFE is stable across all of these; Callaway--Sant'Anna stays near 0.03 and is never significant.
\end{itemize}
\end{frame}

\begin{frame}{What holds up (2): only large dairies show it in TWFE}
\input{generated/headline__robust_spec.tex}
\end{frame}

\begin{frame}{The 2016 data break}
\input{generated/headline__split.tex}
\begin{itemize}\footnotesize\setlength\itemsep{1pt}
  \item CS has the same sign and similar size on both sides (0.023, 0.028); neither is significant.
  \item Pooled TWFE (0.104) exceeds either side's (0.034, 0.062), which fits cross-cohort comparisons; not tested directly. The windows also hold different cohorts.
\end{itemize}
\end{frame}

\begin{frame}{Spatial dependence in the residuals}
\input{generated/headline__gp.tex}
\begin{itemize}\scriptsize\setlength\itemsep{1pt}
  \item Poor MH Days and crime: short-range, mostly noise; Conley errors to 500 km match the clustered ones.
  \item Frequent Mental Distress: long-range, probably from how the CDC series is built.
\end{itemize}
\end{frame}

\begin{frame}{Slaughterhouses: every design}
\input{generated/headline__fsis.tex}
\end{frame}

\begin{frame}{Slaughterhouses: pre-trends and the closure control}
\figcols{generated/assets/headline_fsis_event.pdf}{
\begin{itemize}\setlength\itemsep{4pt}
  \item Entering counties start \emph{higher} and converge down; the negative post estimates track that.
  \item Closure control: the entry estimate stays near $-0.05$; counties after a closure do not have more poor mental health days.
  \item Dropping closure counties flips the sign (+0.124) but fails pre-trends; many-vs-none fails its lead test.
\end{itemize}}
\end{frame}

% =====================================================================
\section{Main takeaway}

\begin{frame}{Scorecard: the claim against its own tests}
\input{generated/headline__scorecard.tex}
\end{frame}

\begin{frame}{What we can claim}
\begin{itemize}\setlength\itemsep{7pt}
  \item \textbf{Association.} Under TWFE, a first expansion of large dairies goes with about 0.10 more poor mental health days per month; that is robust to controls, region-by-year effects, spatial errors, and specific to large dairies.
  \item \textbf{Not yet causal.} The estimators built for staggered timing give about a third of that and are not significant; half of the Callaway--Sant'Anna pre-periods are not zero; the estimate does not survive correction across tests.
  \item \textbf{Slaughterhouses.} No causal pattern: entry is wrong-signed and follows a pre-trend, exit is not estimable, closures are not the explanation.
  \item \textbf{What would change this.} Finer timing of dairy expansions than the five-year census; an outcome measured the same way across years; worker-level exposure data.
\end{itemize}
\end{frame}

\begin{frame}{Decisions for the advisor}
\begin{enumerate}\setlength\itemsep{8pt}
  \item Is ``suggestive association'' an acceptable framing for the headline, or should the paper be descriptive?
  \item Is it worth getting finer dairy timing (for example state permit data) to test the coarse-timing explanation of the pre-period estimates?
  \item Do we pursue a worker-exposure design? It needs employment-share data we do not have.
  \item Do we drop or reframe the slaughterhouse arm, given no design passes its pre-trend rule?
\end{enumerate}
\end{frame}

% =====================================================================
\section{Appendix}

\begin{frame}{Appendix: estimators, first entry}
\fig[0.55]{generated/assets/forest_triangulation_first_entry__main.pdf}
\src{Same layout as the first-expansion frame; Callaway--Sant'Anna 0.026 for Poor MH Days. Source: script4d, 2026-10-24.}
\end{frame}

\begin{frame}{Appendix: pre-period coefficients, first expansion and first entry}
\fig[0.55]{generated/assets/pretrends_leads__main.pdf}
\end{frame}

\begin{frame}{Appendix: TWFE across the four treatments}
\fig[0.55]{generated/assets/forest_a2_treatment_grid__main.pdf}
\end{frame}

\begin{frame}{Appendix: dropping nearby comparison counties, first expansion}
\fig[0.55]{generated/assets/forest_donut_first_expansion__main.pdf}
\end{frame}

\begin{frame}{Appendix: dairy, large operations against all sizes}
\fig[0.55]{generated/assets/forest_species_dairy__main.pdf}
\src{Beef, hogs and chickens are in the species results (script4a, 2026-10-25); the robustness table shows them next to dairy.}
\end{frame}

\begin{frame}{Appendix: crime outcomes, TWFE event studies}
\fig[0.55]{generated/assets/event_study_twfe__backup.pdf}
\src{Arrests, not offences. Pre-trend tests fail for the assault family and the NIBRS total (script4b leads); not part of the headline.}
\end{frame}

\begin{frame}{Appendix: slaughterhouses on their own, plant count and plant entry}
\fig[0.55]{generated/assets/forest_slaughterhouse__main.pdf}
\end{frame}

\end{document}
